% Segunda lectura — Bloque geométrico: transformaciones proyectivas
\section{Bloque geométrico: transformaciones proyectivas}\label{cap:12}\label{sec:bloque-proyectivas}

{\small\vspace{1mm}\linkUnidad{05_proyectivas}{Simulación interactiva: Unidad 7 — Proyectivas: la foto del dron y el pentágono del profesor}\par}
{\small\vspace{1mm}\linkAnimacion{02_homogeneas_pentagono_3d}{Animación: el pentágono en coordenadas homogéneas (3D)}\par}


\section{Homografías: forma algebraica}\label{sec:homografias}

\textbf{Desarrollo.} Una \emph{transformación proyectiva} u \emph{homografía} del plano es la
aplicación
\[
  \boxed{p\longmapsto p'\sim Hp},
\]
con \(H\) una matriz real \(3\times3\) invertible que actúa sobre las coordenadas
homogéneas \(p=(x,y,1)^T\). El símbolo \(\sim\) recuerda que el resultado se lee \emph{a menos
de escala}: se normaliza dividiendo por la tercera coordenada. Como \(H\) y \(\lambda H\)
definen la misma transformación, una homografía tiene \(8\) grados de libertad. Se compone
multiplicando matrices con la convención habitual: \(H_2H_1\) significa primero \(H_1\) y
después \(H_2\).

\begin{resultado}[Qué conserva una homografía]
\begin{itemize}
  \item Rectas en rectas, e incidencia: si \(\ell\cdot p=0\) entonces \(\ell'\cdot p'=0\)
  con \(\ell'=H^{-T}\ell\). Los puntos se transforman con \(H\), las rectas con
  \(H^{-T}\).
  \item La \emph{razón doble} de cuatro puntos alineados, el invariante fundamental de la
  geometría proyectiva (no se desarrolla en este apunte).
  \item No conserva distancias, ángulos, áreas, paralelismo ni razones simples: el punto medio
  de un segmento puede ir a cualquier punto interior de la imagen.
\end{itemize}
\end{resultado}

\paragraph{La recta que se va al infinito.}
Sea \(h_3=(h_{31},h_{32},h_{33})\) la tercera fila de \(H\). La tercera coordenada de la
imagen es \(W=h_3\cdot p\). Los puntos de la recta \(h_{31}x+h_{32}y+h_{33}=0\) tienen
\(W=0\): se van al infinito. Si esa recta \emph{no toca} una figura acotada, la imagen es una
figura acotada con las mismas incidencias. Si la \emph{atraviesa}, los puntos de un lado
tienen \(W>0\) y los del otro \(W<0\), y la imagen se parte: un segmento que cruza la recta se
transforma en dos semirrectas unidas ``por el infinito'', como un objeto que está en parte
delante y en parte detrás de la cámara.

\begin{atencion}[Las afinidades vistas desde acá]
Una afinidad tiene \(h_3=(0,0,1)\): la recta que se va al infinito es la propia recta ideal.
Ningún punto ordinario se pierde, las paralelas siguen paralelas y las figuras nunca se
parten. Recíprocamente, toda homografía que deja fija la recta ideal es una afinidad.
\end{atencion}

\section{Cuatro puntos determinan la homografía}\label{sec:cuatro-puntos}

\paragraph{Conteo.} Cada correspondencia \(p_i\mapsto p_i'\) aporta dos ecuaciones
(\(X/W=x_i'\), \(Y/W=y_i'\)). Cuatro correspondencias dan ocho ecuaciones para ocho
incógnitas. La condición es que ni entre los cuatro puntos de partida ni entre los cuatro de
llegada haya tres alineados (con el determinante se verifica).

\paragraph{Construcción.} En lugar de resolver el sistema \(8\times8\) se pasa por un
\emph{cuadrilátero patrón}:
\[
  e_1=\begin{pmatrix}1\\0\\0\end{pmatrix},\quad
  e_2=\begin{pmatrix}0\\1\\0\end{pmatrix},\quad
  e_3=\begin{pmatrix}0\\0\\1\end{pmatrix},\quad
  e_4=\begin{pmatrix}1\\1\\1\end{pmatrix}=e_1+e_2+e_3 .
\]
Dados cuatro puntos \(p_1,p_2,p_3,p_4\) sin tres alineados, se busca una matriz \(K\) con
\(Ke_i\sim p_i\) para \(i=1,\dots,4\):
\begin{enumerate}
  \item Se forma \(M=[\,p_1\ p_2\ p_3\,]\) con los tres primeros puntos como columnas. Como
  \(Me_i\) es la \(i\)-ésima columna, ya se cumple \(Me_i=p_i\) para \(i\le3\). Pero
  \(Me_4=p_1+p_2+p_3\), que en general no es \(p_4\).
  \item Como cada columna representa su punto \emph{a menos de escala}, se la puede
  multiplicar por un número sin cambiar a dónde va \(e_i\). Se eligen \(\lambda_1,\lambda_2,
  \lambda_3\) tales que
  \[
    \lambda_1p_1+\lambda_2p_2+\lambda_3p_3=p_4
    \qquad\Longleftrightarrow\qquad
    \lambda=M^{-1}p_4 .
  \]
  \item La matriz buscada es \(K=[\,\lambda_1p_1\ \ \lambda_2p_2\ \ \lambda_3p_3\,]\):
  las tres primeras columnas siguen representando \(p_1,p_2,p_3\) y su suma es \(p_4\).
\end{enumerate}

\begin{resultado}[Homografía que lleva un cuadrilátero sobre otro]
Sean \(K_1\) la matriz del patrón para los puntos de partida \(p_1,\dots,p_4\) y \(K_2\) la
del patrón para los puntos de llegada \(p_1',\dots,p_4'\). Entonces
\[
  \boxed{H=K_2K_1^{-1}}
\]
es la única homografía (a menos de escala) con \(Hp_i\sim p_i'\), \(i=1,\dots,4\).
Se lee de derecha a izquierda: \(K_1^{-1}\) lleva el cuadrilátero de partida al patrón, y
\(K_2\) lleva el patrón al cuadrilátero de llegada.
\end{resultado}

Es el mismo esquema del ejercicio 6 de afinidades (\(A=P'P^{-1}\)), con un punto más para
absorber la perspectiva. Los \(\lambda_i\), invisibles punto a punto, son precisamente los
grados de libertad que faltaban.

\section{Puntos fijos y rectas fijas de una homografía}\label{sec:homografia-fijos}

\textbf{Desarrollo.} Acá se cierra la cadena del bloque anterior. \(p\) es fijo si
\(Hp\sim p\), es decir si existe \(\lambda\neq0\) con \(Hp=\lambda p\): \emph{los puntos
fijos son los autovectores reales de \(H\)}. Se calculan como se explicó: raíces reales de
\(\det(H-\lambda I)=0\), y para cada una el producto cruz de dos filas de \(H-\lambda I\).

\begin{resultado}[Existencia de punto fijo]
Toda homografía del plano tiene al menos un punto fijo, porque el polinomio característico
es de grado \(3\). Es ordinario si el autovector tiene tercera coordenada no nula, e ideal
(una dirección fija) si la tercera coordenada es cero. Una raíz doble con producto cruz de
filas nulo significa toda una recta de puntos fijos.
\end{resultado}

\paragraph{Rectas fijas.} Como las rectas se transforman con \(H^{-T}\), una recta es fija
si es autovector de \(H^{T}\); se encuentra con el producto cruz de dos filas de
\(H^{T}-\lambda I\) (o sea, de dos columnas de \(H-\lambda I\)). Una recta fija de autovalor
\(\mu\) pasa por todos los puntos fijos de autovalor \(\lambda\neq\mu\).

\paragraph{Clasificación por autovalores.}
\begin{itemize}
  \item \textbf{Tres autovalores reales distintos:} tres puntos fijos no alineados y tres
  rectas fijas, los lados de ese triángulo.
  \item \textbf{Un real y dos complejos conjugados:} un único punto fijo real y una única
  recta fija real que no pasa por él. En coordenadas adecuadas es una rotación alrededor del
  punto fijo, con la recta fija en el papel de recta ideal.
  \item \textbf{Un autovalor doble con producto cruz de filas nulo:} una recta entera de
  puntos fijos (el \emph{eje}) y un punto fijo aislado (el \emph{centro}) con el otro
  autovalor. Es una \emph{homología}. Si los dos autovalores son opuestos, \(H^2\sim I\): la
  homología es \emph{armónica}, una involución, como una reflexión.
\end{itemize}

\section{Jerarquía de transformaciones del plano}\label{sec:jerarquia}

\begin{center}
\begin{tabularx}{\textwidth}{>{\bfseries}p{24mm} p{44mm} c c X}
\toprule
Familia & Forma de \(H\) & g.d.l. & Puntos & Conserva\\
\midrule
Isometrías & \(\begin{pmatrix}R&t\\0&1\end{pmatrix}\), \(R^TR=I\) & 3 & 2 &
distancias, ángulos, áreas\\
Similaridades & \(\begin{pmatrix}kR&t\\0&1\end{pmatrix}\) & 4 & 2 &
ángulos, razones de longitudes\\
Afinidades & \(\begin{pmatrix}A&t\\0&1\end{pmatrix}\), \(\det A\neq0\) & 6 & 3 &
paralelismo, razón simple, razón de áreas\\
Proyectivas & \(H\) invertible cualquiera & 8 & 4 &
rectas, incidencia, razón doble\\
\bottomrule
\end{tabularx}
\end{center}
``Puntos'' es la cantidad de correspondencias \(p_i\mapsto p_i'\) que determinan la
transformación. Cada familia contiene a la anterior y conserva menos.

\section{El ejercicio del pentágono regular}

\begin{enunciado}[Dictado, clase del 18 de septiembre]
Tomar un pentágono \emph{regular} \(ABCDE\). Una transformación proyectiva queda definida
por cuatro puntos, así que se prescriben cuatro vértices y sus imágenes (otros cuatro
vértices) y se pregunta \emph{a dónde cae el quinto}. Hay tres casos, y sus combinaciones
con las rotaciones del pentágono:
\begin{enumerate}[label=(\alph*)]
  \item no intercambiar ningún vértice: sólo ``correr'' el pentágono;
  \item intercambiar dos vértices adyacentes;
  \item intercambiar dos vértices no adyacentes.
\end{enumerate}
Para cada transformación, determinar si tiene punto fijo (un autovector de la matriz).
Entrega: un único PDF en el campus (o por correo), con figuras y observaciones; se
comparten avances en el grupo.
\end{enunciado}

\textbf{Desarrollo.} Se toma el pentágono regular inscripto en la circunferencia unidad,
centrado en el origen y con un vértice arriba. El docente sugirió centrarlo en el origen:
en el caso más simple el punto fijo es el centro, cuyo autovector es \((0,0,1)^T\).
\begin{gather*}
  A=(0,\,1),\qquad B=(-0.9511,\,0.3090),\qquad C=(-0.5878,\,-0.8090),\\
  D=(0.5878,\,-0.8090),\qquad E=(0.9511,\,0.3090).
\end{gather*}
Las coordenadas son irracionales (\(\cos72^\circ=\frac{\sqrt5-1}{4}\)), así que las
matrices se calculan con cuatro decimales y se verifican en NumPy (apéndice~\ref{ap:pentagono}).
En los tres casos se usa \(H=K_2K_1^{-1}\) con \(K_1\) el patrón de \(ABCD\), y se
normaliza \(h_{33}=1\).

\begin{center}
\begin{tikzpicture}[scale=1.9]
  \draw[linegray!40,very thin] (-1.6,-1.4) grid[step=.5] (1.6,1.3);
  \draw[axisline,-{Latex[length=2.2mm]}] (-1.7,0)--(1.7,0) node[right] {$x$};
  \draw[axisline,-{Latex[length=2.2mm]}] (0,-1.45)--(0,1.35) node[above] {$y$};
  \draw[linegray,densely dotted] (0,0) circle (1);
  \draw[teal,line width=1.05pt,fill=teallight!45] (0,1)--(-0.9511,0.3090)--(-0.5878,-0.8090)--(0.5878,-0.8090)--(0.9511,0.3090)--cycle;
  \foreach \p/\n/\pos in {(0,1)/A/above,(-0.9511,0.3090)/B/left,(-0.5878,-0.8090)/C/{below left},(0.5878,-0.8090)/D/{below right},(0.9511,0.3090)/E/right}
    \fill[teal] \p circle (1.2pt) node[\pos,font=\small,text=tealdark] {$\n$};
  \fill[ink] (0,0) circle (1.2pt) node[below right,font=\scriptsize] {$O$};
  % caso (c): imagen C B A D E'
  \draw[rust,line width=1pt] (-0.5878,-0.8090)--(-0.9511,0.3090)--(0,1)--(0.5878,-0.8090)--(0,-1.2361)--cycle;
  \fill[rust] (0,-1.2361) circle (1.4pt) node[right=2pt,font=\small,text=rust,labelback] {$E'_c=(0,-1.2361)$};
  % eje BD
  \draw[dashed,tealdark,line width=.8pt] ($(-0.9511,0.3090)!-0.35!(0.5878,-0.8090)$)--($(-0.9511,0.3090)!1.35!(0.5878,-0.8090)$) node[right,font=\scriptsize,text=tealdark] {eje $BD$};
  % recta AC (contiene al centro)
  \draw[dotted,gold!70!black,line width=.8pt] ($(0,1)!-0.3!(-0.5878,-0.8090)$)--($(0,1)!1.3!(-0.5878,-0.8090)$);
  \node[gold!60!black,font=\scriptsize,labelback,anchor=east] at (-0.55,-1.3) {hacia $O_c=(-1.54,-3.74)$};
\end{tikzpicture}
\end{center}

\paragraph{(a) Sin intercambio: \(A\mapsto B,\ B\mapsto C,\ C\mapsto D,\ D\mapsto E\).}
Los cuatro pares son compatibles con la rotación de \(72^\circ\) alrededor de \(O\), que es
una isometría y por lo tanto una homografía. Como la homografía es única, \emph{es} esa
rotación:
\[
  H_a=\begin{pmatrix}0.3090&-0.9511&0\\0.9511&0.3090&0\\0&0&1\end{pmatrix},
  \qquad H_aE=A .
\]
El quinto vértice cae en \(A\): el pentágono regular sí ``cierra''. Autovalores \(1\) y
\(e^{\pm i\,72^\circ}\); único punto fijo real \(O=(0,0)\), obtenido como producto cruz de
las dos primeras filas de \(H_a-I\):
\[
  (-0.691,\,-0.951,\,0)\times(0.951,\,-0.691,\,0)=(0,\,0,\,1.382)\sim(0,0,1)^T .
\]
La recta ideal es fija (es afín) y ninguna otra.

\paragraph{(b) Intercambio de adyacentes: \(A\leftrightarrow B\), \(C\) y \(D\) quietos.}
\[
  H_b=\begin{pmatrix}-5.8541&-4.9798&-4.0287\\0&-1&3.9271\\0&8.4721&1\end{pmatrix},
  \qquad E\mapsto E'_b=(-3.0777,\ 1.0000).
\]
Autovalores \(5.8541\) (simple) y \(-5.8541\) (doble). Para el doble, el producto cruz de
cualquier par de filas de \(H_b+5.8541I\) es nulo: hay toda una recta de puntos fijos, que es
la recta \(CD\) (\(y=-0.8090\)), lógico porque \(C\) y \(D\) quedan quietos. El autovalor
simple da el centro, \(O_b=(-0.5878,\,0.5729)\), que está sobre la recta \(AB\) (la recta
\(AB\) va a sí misma porque sus extremos se intercambian). Como los autovalores son
opuestos, \(H_b^2\sim I\): es una \emph{homología armónica}, un ``espejo proyectivo'' de eje
\(CD\) y centro \(O_b\). La recta que se va al infinito es \(y=-0.1180\), que
\emph{atraviesa} el pentágono (\(W\) vale \(9.47,\ 3.62,\ -5.85,\ -5.85,\ 3.62\) en
\(A,\dots,E\)): los lados \(BC\) y \(DE\) la cruzan y la imagen del pentágono se parte en
dos semirrectas por lado.

\paragraph{(c) Intercambio de no adyacentes: \(A\leftrightarrow C\), \(B\) y \(D\) quietos.}
\[
  H_c=\begin{pmatrix}0.4271&-0.5878&-0.2245\\-1.0368&-0.5729&-0.5451\\0.2775&0.3820&1\end{pmatrix},
  \qquad E\mapsto E'_c=(0,\ -1.2361).
\]
Autovalores \(0.8541\) (doble) y \(-0.8541\). Eje de puntos fijos: la recta \(BD\)
(\(0.7265x+y+0.3820=0\)); centro \(O_c=(-1.5388,\,-3.7361)\), sobre la prolongación de
\(AC\), obtenido como producto cruz de las dos primeras filas de \(H_c+0.8541I\):
\((1.2812,-0.5878,-0.2245)\times(-1.0368,0.2812,-0.5451)=(0.3835,\,0.9311,\,-0.2492)\sim
(-1.5388,-3.7361,1)\). Otra vez \(H_c^2\sim I\): homología armónica de eje \(BD\). El
quinto vértice cae en \(E'_c=(0,-1.2361)\), sobre el eje de simetría, y como la recta que se
va al infinito (\(0.2775x+0.3820y+1=0\)) no toca el pentágono (\(W>0\) en los cinco
vértices), la imagen es el pentágono \(CBADE'_c\) de la figura, no convexo pero entero.

\paragraph{Combinaciones con rotaciones.} Cualquier otra correspondencia de cuatro vértices
sobre cuatro vértices es una de las tres anteriores seguida de una rotación \(H_a^k\), o una
reflexión del pentágono. Por ejemplo \(A\mapsto E,\ B\mapsto D,\ C\mapsto C,\ D\mapsto B\)
es la reflexión respecto del eje \(OC\) (afín, \(E\mapsto A\), todo el eje fijo y la
dirección perpendicular como punto fijo ideal). Componer (b) con la rotación da un único punto
fijo real y dos autovalores complejos; componer (c) con la rotación da tres autovalores
reales distintos y un triángulo invariante. Es un buen ejercicio de exploración con el
script del apéndice.

\begin{resultado}[Respuestas a la consigna]
\begin{itemize}
  \item Hay \(600\) correspondencias de cuatro vértices ordenados en cuatro vértices, todas
  bien definidas (el pentágono es convexo); a menos de las \(10\) simetrías del pentágono se
  reducen a los casos (a), (b), (c) y sus composiciones.
  \item (a) es la rotación; el quinto vértice cae en un vértice y el punto fijo es el centro.
  \item (b) y (c) son homologías armónicas: una recta entera de puntos fijos (el lado cuyos
  extremos quedan quietos) más un centro; el quinto vértice cae fuera del pentágono, en
  \((-3.08,1)\) y \((0,-1.24)\) respectivamente. En (b) la imagen se parte; en (c) no.
  \item Siempre hay punto fijo: es el autovector real de \(H\), y se calcula con el producto
  cruz de dos filas de \(H-\lambda I\).
\end{itemize}
\end{resultado}

\section{Ejercicios de consolidación: proyectivas}

\begin{enumerate}
  \item Verifique con el determinante que \(A,B,C\) del pentágono regular no están alineados
  y que \(A,O,C\) tampoco (\(O\) el centro).
  \item Calcule la recta \(BD\) como producto vectorial de las ternas de \(B\) y \(D\).
  Compruebe que el centro \(O_c=(-1.5388,-3.7361)\) \emph{no} está en ella.
  \item Con \(H_c\), calcule la imagen del centro \(O\) y compruebe que no es \(O\): en
  (c) el centro del pentágono no queda fijo.
  \item Calcule \(H_c^2\) y normalice: debe dar la identidad. Interprete.
  \item Construya la matriz \(K\) del patrón para el cuadrado \((0,0),(1,0),(1,1),(0,1)\),
  en ese orden. ¿Es afín la homografía que lleva el patrón al cuadrado? Explique por qué.
\end{enumerate}

\subsection{Respuestas para autocorrección}
\begin{enumerate}
  \item \(\det[A\ B\ C]=1.3143\neq0\); \(\det[A\ O\ C]=-0.5878\neq0\).
  \item \(B\times D=(0.3090+0.8090,\ 0.5878+0.9511,\ 0.7694-0.1816)\cdot\) (a menos de
  escala) \(\sim(0.7265,\,1,\,0.3820)\); en \(O_c\): \(-1.118-3.736+0.382=-4.47\neq0\).
  \item \(H_c(0,0,1)^T=(-0.2245,-0.5451,1)^T\): \(O\mapsto(-0.22,-0.55)\).
  \item \(H_c^2=0.7295\,I\): a menos de escala, la identidad; \(H_c\) es una involución.
  \item \(M=\begin{pmatrix}0&1&1\\0&0&1\\1&1&1\end{pmatrix}\), \(\lambda=(1,-1,1)\),
  \(K=\begin{pmatrix}0&-1&1\\0&0&1\\1&-1&1\end{pmatrix}\). No es afín (última fila
  \((1,-1,1)\)): el patrón contiene los puntos ideales \(e_1,e_2\) y los manda a puntos
  propios, cosa que una afinidad nunca hace.
\end{enumerate}

